\documentclass[10pt,1p,authoryear]{elsarticle}
\usepackage[utf8]{inputenc}

\setlength{\tabcolsep}{1.3pt}
\setlength{\columnsep}{0.4cm}
\usepackage{acronym}
\usepackage{amsmath}
\usepackage{amssymb}
\usepackage{array}
\usepackage{bm}
\usepackage{booktabs}
\usepackage[aboveskip=1pt,labelfont=bf,labelsep=period,justification=raggedright,singlelinecheck=off, font=small]{caption}
\usepackage{diagbox}
\usepackage{float}
\usepackage[T1]{fontenc}
% \usepackage[a4paper,inner=2cm,outer=2cm,top=2.5cm,bottom=2.5cm]{geometry}
\usepackage{graphicx}% Include figure files
\usepackage{here}
\usepackage{hyperref}
%\usepackage{lipsum}
\usepackage{lineno}
\usepackage{setspace}
\usepackage{mathtools}
\usepackage{mathrsfs}
\usepackage{multirow}
\usepackage{newfloat}
\usepackage{onimage}
\usepackage{siunitx}
\usepackage{subcaption}
\usepackage{tabularx}
\usepackage{tikz}
%\usepackage[dvipsnames]{xcolor}
\usepackage{xpatch}
\usepackage{orcidlink}

% Ensure that subcaptions have capital letter
\renewcommand{\thesubfigure}{\Alph{subfigure}}

% Must be loaded after tikz
\usepackage{ctable}

% Define new column type with multirow
\newcolumntype{M}[1]{>{\centering\arraybackslash}m{#1}}

\bibliographystyle{elsarticle-harv}

\restylefloat{table}

% Remove indent across the whole document
\setlength\parindent{0pt}

% Define equation numbering
\numberwithin{equation}{section}


\newcommand{\mbx}{\mathbf{x}}
\newcommand{\mbm}{\mathbf{m}}
\newcommand{\mbp}{\mathbf{p}}
\newcommand{\mbu}{\mathbf{u}}
\newcommand{\mbF}{\mathbf{F}}
\newcommand{\mbH}{\mathbf{H}}
\newcommand{\cmt}[1]{\textcolor{red} {[#1]}}
\newcommand{\etal}{{\textit{et al. }}}

\newacro{ode}[ODE]{Ordinary Differential Equation}
\newacro{pde}[PDE]{Partial Differential Equation}
\newacro{glv}[gLV]{generalized Lotka-Volterra}
\newacro{abm}[ABM]{Agent-Based Model}


% Define new style for json listings
% Copied from https://www.latex4technics.com/?note=3QTU
\colorlet{punct}{red!60!black}
\definecolor{background}{HTML}{EEEEEE}
\definecolor{delim}{RGB}{20,105,176}
\colorlet{numb}{magenta!60!black}

\title{The Pool Model: A Mechanistic ODE Framework for Bacterial Consortia in Food Science}
\author[1]{Polina Gaindrik\orcidlink{0000-0003-0928-3951}}
\author[1]{Jonas Pleyer\orcidlink{0009-0001-0613-7978}}
\author[2]{Matthias Brunner}
\author[3]{Fernando Pérez Rodríguez}
\author[1]{Christian Fleck\orcidlink{0000-0002-6371-4495}\corref{cor1}}
\ead{christian.fleck@fdm.uni-freiburg.de}

\cortext[cor1]{Corresponding Author\\ \textit{E-mail address:} christian.fleck@fdm.uni-freiburg.de}
\affiliation[1]{
    organization={Freiburg Center for Data Analysis and Modeling and AI},
    addressline={Ernst-Zermelo Str. 1},
    city={Freiburg im Breisgau},
    postcode={79104},
    country={Germany},
}
\affiliation[2]{
    organization={tsenso GmbH},
    addressline={Nöllenstr. 32},
    city={Stuttgart},
    postcode={70195},
    country={Germany},
}
\affiliation[3]{
    organization={Department of Food Science and Technology, UIC Zoonosis y Enfermedades Emergentes ENZOEM, ceiA3, Universidad de Córdoba},
    addressline={Campus Rabanales s/n. Darwin building- Annex},
    city={Córdoba},
    postcode={14014},
    country={Spain},
}
\date{
    \footnotesize
    \textsuperscript{\textbf{1}}
}

% Force double spacing in abstract (elsarticle overrides \doublespacing there)
\makeatletter
\let\els@abstract\abstract
\renewcommand{\abstract}{\els@abstract\setstretch{2}}
\makeatother

\begin{document}
\doublespacing
\linenumbers
%\maketitle
%
% ================================================================================================% ================================================================================================
\section{Case Study}

%@article{ferrer-bustinsAntilisterialEffect2024,
%	title = {The Antilisterial Effect of Latilactobacillus sakei {CTC}494 in Relation to Dry Fermented Sausage Ingredients and Temperature in Meat Simulation Media},
%	volume = {10},
%	rights = {http://creativecommons.org/licenses/by/3.0/},
%	issn = {2311-5637},
%	url = {http://dx.doi.org/10.3390/fermentation10060326},
%	doi = {10.3390/fermentation10060326},
%	pages = {326},
%	number = {6},
%	journaltitle = {Fermentation},
%	publisher = {Multidisciplinary Digital Publishing Institute},
%	author = {Ferrer-Bustins, N鷕ia and Costa, Jean Carlos Correia Peres and P閞ez-Rodr韌uez, Fernando and Mart韓, Bel閚 and Bover-Cid, Sara and Jofr�, Anna},
%	urldate = {2026-06-18},
%	date = {2024-06},
%	langid = {english},
%   year={2024}
%}

%@article{leroyModelingBacteriocin2005,
%	title = {Modeling Bacteriocin Resistance and Inactivation of Listeria innocua {LMG} 13568 by Lactobacillus sakei {CTC} 494 under Sausage Fermentation Conditions},
%	volume = {71},
%	issn = {0099-2240},
%	url = {http://dx.doi.org/10.1128/AEM.71.11.7567-7570.2005},
%	doi = {10.1128/AEM.71.11.7567-7570.2005},
%	pages = {7567--7570},
%	number = {11},
%	journaltitle = {Applied and Environmental Microbiology},
%	shortjournal = {Appl Environ Microbiol},
%	author = {Leroy, Frédéric and Lievens, Kristoff and De Vuyst, Luc},
%	date = {2005-11},
%    year={2005}
%}

To demonstrate the efficacy of the proposed approach using real experimental data, we consider a two-species system representative of a food-based scenario.
As described by~\citet{ferrer-bustinsAntilisterialEffect2024} the dynamics between \textit{Listeria monocytogenes} CTC1034 (Lm-CTC1034), a pathogen common in dry fermented sausages, and two strains of  \textit{Latilactobacillus sakei} were studied using a central composite design in meat simulation media.
These strains included CTC494 (Ls-CTC494), known for its production of the bacteriocin sakasin K, which exhibits strong antimicrobial activity against \textit{Listeria spp.}, and 23K (Ls-23K), which does not produce bacteriocins.
Growth curves were measured for each species (Lm-CTC1034, Ls-CTC494, Ls-23K) in monoculture, as well as in cocultures combining \textit{L. monocytogenes} with each \textit{L. sakei} strain.
Furthermore, time-series data were collected for bacteriocin activity, pH values, and the concentration of lactic acid produced by both species.
Accordingly, we propose the two-species model for \textit{L. sakei} and \textit{L. monocytogenes}.
Based on the experimental data, we assume that the bacteria had already passed the lag phase and entered the growth phase.
So the transition from the lag pool is omitted from our model. 
Both species compete for a common resource pool $\tilde{R}$ with capacity $N_t$.
Moreover, the toxin produced during the growth of Ls-CTC494 at rate $\kappa_T^{\text{Ls } \delta}$ induces the death of Lm-CTC1034 cells via a Hill-function dependence with a rate $\omega'_\text{Lm}$.
As it was observed in works with \textit{L. innocua}, e.g., by~\citet{leroyModelingBacteriocin2005}, small number of \textit{Listeria spp.} cells is resistant to sakasin K.
Given that we divide the Lm-CTC1034 population into two groups: sensitive ($x^\text{sens}_\text{Lm}$) and resistant ($x^\text{res}_\text{Lm}$) to it.
For simplicity, we assume identical growth dynamics ($\mu^\text{res}_\text{Lm} = \mu^\text{res}_\text{Lm} = \mu_\text{Lm}$) for both groups, with the toxin-induced death term applying only to the sensitive population.
Both species continuously produce lactic acid at rates $\kappa_{LA}^{\text{Ls } \delta}$ and $\kappa_{LA}^\text{Lm}$,  while \textit{L. sakei} also exhibits growth-associated production at a rate $\kappa_{LA/G}^{\text{Ls } \delta}$.
Based on experimental observations of lactic acid concentration, we assume that metabolic differences between the Ls-23K and Ls-CTC494 strains lead to different production rates.
Thus, the index $\delta$ is used to differentiate parameters between the two \textit{L. sakei} strains. 
Note that the toxin production rate $\kappa_T^{\text{Ls } \delta}$  also depends on the strain of \textit{L. sakei}: for Ls-23K it equals zero and for Ls-CTC494 it is positive value.
The growth rates for two \textit{L. sakei} strains are assumed to be equal.
\begin{align}
    \begin{split}
        \dot x^\delta_\text{Ls} &= \mu_\text{Ls} \tilde{R} x^\delta_\text{Ls} \\%- \omega_\text{Ls} x_\text{Ls} \\
        \dot x^\text{sens}_\text{Lm} &= \mu_\text{Lm} \tilde{R} x^\text{sens}_\text{Lm} - \omega'_\text{Lm} \frac{\mathcal{T}^n_\text{Ls}}{\mathcal{T}^n_\text{Ls} + K^n} x^\text{sens}_\text{Lm} \\ %  - \omega_\text{Lm} x_\text{Lm}
        \dot x^\text{res}_\text{Lm} &= \mu_\text{Lm} \tilde{R} x^\text{res}_\text{Lm} \\
        \dot {\tilde{R}} &= - \frac{1}{N_t} (\mu_\text{Ls} \tilde{R} x^\delta_\text{Ls} + \mu_\text{Lm} \tilde{R} x^\text{sens}_\text{Lm} + \mu_\text{Lm} \tilde{R} x^\text{res}_\text{Lm}) \\
        \dot{\mathcal{T}_\text{Ls}} &= \kappa_T^{\text{Ls } \delta} \tilde{R} x^\delta_\text{Ls} \\
        \dot {LA} &= (\kappa_{LA}^{\text{Ls } \delta} + \kappa_{LA/G}^{\text{Ls } \delta} \tilde{R}) x^\delta_\text{Ls} + \kappa_{LA}^\text{Lm} (x^\text{sens}_\text{Lm} + x^\text{res}_\text{Lm})%\\
        %\dot {pH} &= q(t) \dot {LA}.
        \label{eq:case_study_ode}
    \end{split}
\end{align}
Furthermore, the accumulation of lactic acid precipitates a decrease in the pH of the microenvironment, which in turn alters the growth dynamics.
To account for this effect, we incorporate a term from the gamma model proposed by~\citet{zwietering_decision_1992}:
\begin{equation}
    \mu (pH) = \mu_\text{opt} \gamma(pH) =  \mu_\text{opt} \frac{pH - pH_\text{min}}{pH_\text{opt} - pH_\text{min}},
    \label{eq:case_study_mu_pH}
\end{equation}
where $\mu_\text{opt}$ is the growth rate at optimum conditions.
Based on the experimental data we estimated parameter of the model (see Fig.~\ref{fig:case_study}). 
%
\begin{figure}[H]
    \begin{center}
    \includegraphics[width=0.8\columnwidth]{Figures-pool_model_real_data_all_count.pdf}\\
    \includegraphics[width=0.8\columnwidth]{Figures-pool_model_real_data_LA_pH.pdf} 
    \caption{
        The model~(\ref{eq:case_study_ode})  is shown with parameters estimated based on experimental data (represented by scatter plots). Dashed lines denote model solutions for monoculture experiments, while solid lines denote coculture experiments.
        (A) Bacterial counts for monocultures of \textit{L. sakei} 23K (Ls-23K, light green), \textit{L. sakei} CTC494 (Ls-CTC494, light blue), and \textit{L. monocytogenes} CTC1034 (Lm-CTC1034, orange).
        Cocultures are shown for Lm-CTC1034 with Ls-23K (Ls-23K in dark green, Lm-CTC1034 in pink) and Lm-CTC1034 with Ls-CTC494 (Ls-CTC494 in dark blue, Lm-CTC1034 in red). The activity of the bacteriocin produced by Ls-CTC494 is indicated by the brown dotted line.
        (B) Lactic acid concentrations for monocultures of Ls-23K (light green), Ls-CTC494 (light blue), and Lm-CTC1034 (orange), as well as cocultures of Lm-CTC1034 with Ls-23K (purple) and Lm-CTC1034 with Ls-CTC494 (yellow). Measured pH values, used as an external condition, are represented by crosses; colors correspond to the respective lactic acid curves.
    }
    \label{fig:case_study}
    \end{center}
\end{figure}

% Add values + units
\begin{table}[H]
    \centering\small
    \begin{tabular}{ccccc}
        \textbf{Parameters} \hspace{5mm}& \textbf{Ls-23K} \hspace{5mm}& \textbf{Ls-CTC494} \hspace{5mm}& \textbf{Lm-CTC1034(sens)} \hspace{5mm}& \textbf{Lm-CTC1034(res)} \\
    \toprule
    \multicolumn{5}{c}{Growth associated parameters}\\
    \midrule
    $\mu_\text{opt}$ &  \multicolumn{2}{c}{ - } & \multicolumn{2}{c}{ - }\\
    $pH_\text{min}$ &  \multicolumn{2}{c}{ - } & \multicolumn{2}{c}{ - }\\
    $pH_\text{opt}$ &  \multicolumn{2}{c}{ - } & \multicolumn{2}{c}{ - }\\
    $N_t$ & \multicolumn{4}{c}{-}\\
    \midrule
    \multicolumn{5}{c}{Lactic acid production}\\
    \midrule
    $\kappa_\text{LA}$ & - & -  &\multicolumn{2}{c}{ - } \\
    $\kappa_\text{LA/G}$ & - & -  &\multicolumn{2}{c}{ 0 } \\
    \midrule
    \multicolumn{5}{c}{Bacteriocin associated parameters}\\
    \midrule
    $\kappa_T^{\text{Ls} \delta}$ & \multicolumn{4}{c}{ - } \\
    $\omega'_\text{Lm}$ & \multicolumn{4}{c}{-}\\
    $K$ & \multicolumn{4}{c}{-}\\
    \midrule
    \multicolumn{5}{c}{Initial conditions}\\
    \midrule
    Monoculture         & - & - & - & - \\
    Coculture (23K)     & - & 0 & - & - \\
    Coculture (CTC494)  & 0 & - & - & - \\
    \end{tabular}
    \caption{
        Model parameters estimated form the experimental data.
     }
    \label{tab:case_study_parameters}
\end{table}

% ================================================================================================
% \printbibliography
%\bibliography{references.bib}

% ================================================================================================
\newpage
\end{document}
